\chapter{Implementation \& AI Evolution}

\section{Introduction}
This chapter delves into the practical execution of the Kemora project, framing its development as a journey of overcoming complex engineering challenges. Tracking its evolution through version control, we explore how abstract designs were translated into concrete code. We place special emphasis on the iterative refinement of the AI prompt engineering pipeline, detailing the bottlenecks encountered and the architectural pivots required to achieve a scalable, low-latency system.

\section{Development Lifecycle and Phases}
An analysis of the full commit history reveals that the project was developed iteratively over seven distinct phases spanning approximately seven months. Rather than a straight path, this lifecycle was defined by continuous testing, profiling, and redesigning in response to real-world constraints.

\subsection{Phase 1: Initial Setup \& Project Scaffold (December 2025)}
The project commenced with repository initialization, establishing the `.gitignore`, and outlining the skeletal structure of the .NET Clean Architecture \cite{clean_architecture} solution. The two December 2025 commits scaffolded the Visual Studio solution and the empty project shells for the Domain, Application, Infrastructure, and API layers. No application functionality was implemented in this phase; the committed code at this point consisted entirely of template-generated files.

\subsection{Phase 2: Database, Auth \& Core API (February 2026)}
The backend architecture matured significantly during this period. The core API controllers were implemented alongside their respective Data Transfer Objects (DTOs) and application services. Security was a primary focus: JWT Authentication was integrated, refresh tokens were added, SMTP email (for email confirmation and password reset) was wired in, and secure secret management was established via .NET User Secrets and environment variables. Initial EF Core migrations and a local database seeder were implemented to populate governorates and base place data.

\subsection{Phase 3: Frontend Architecture (March 2026)}
The Flutter application (\texttt{kemora\_app}) was initialized. The frontend team established the state management architecture and began wiring up the authentication screens to the live backend. By mid-March, the preliminary interfaces for the AI Trip Planning feature were merged into the main branch.

\subsection{Phase 4: Advanced AI Pipeline \& Infrastructure (April - May 2026)}
This phase marked the most significant architectural shift in the backend. On May 2, 2026, the entire AI infrastructure was overhauled. We introduced the \texttt{OpenRouterAiService} \cite{openrouter} to dynamically orchestrate multiple Large Language Models, escaping vendor lock-in. However, we quickly discovered that relying on external LLMs introduced unacceptable latency and cost overheads. To mitigate this, we engineered a dual-layer \texttt{PrecomputedTripPlan} cache, storing successful generations in memory and persisting them to the database to instantly fulfill future identical queries.

\subsection{Phase 5: Testing, Assets \& UI/UX Polish (June 2026)}
Mid-June saw massive refactoring on the client side. A comprehensive End-to-End (E2E) integration test suite was introduced to ensure application stability. Cloudinary was integrated for robust, CDN-backed image hosting. Concurrently, the frontend received a major visual overhaul—dubbed the "Modern Heritage" design system—which refined typography, spatial layouts, and color theory.

\subsection{Phase 6: Performance Optimization (Late June 2026)}
The final weeks were dedicated to significant performance optimization. We solidified our reliance on the meticulously pre-populated local database, entirely removing any dynamic synchronization with external APIs like Google Places to guarantee consistently low latency. The AI pipeline's major bottleneck—fetching images sequentially—was resolved by completely removing external image requests during generation, instead re-injecting image URLs from the pre-loaded local dataset into the AI's final selections.


\section{The Evolution of the AI Pipeline}
The Kemora AI Trip Planner is not a simple passthrough proxy to an LLM. It is a highly orchestrated, multi-stage pipeline designed to enforce geographic constraints and output deterministic data structures.

\subsection{Dynamic Model Selection and Resiliency}
The \texttt{OpenRouterAiService} diverges from traditional AI implementations by eschewing hardcoded models (e.g., \texttt{gpt-4}). Instead, it leverages a dynamic model rotation architecture to utilize free open-source models while maintaining high availability.

\textbf{Model Discovery and Filtering:}
Upon initialization, the service fetches the model directory from the OpenRouter API. The models are rigorously filtered to isolate text-generation models capable of handling extensive RAG context windows (minimum 8192 tokens) while actively blacklisting unreliable, specialized, or paid models that falsely advertise as free. This ensures only high-quality, cost-effective models are used.

\textbf{Intelligent Rotation Strategy:}
The \texttt{SelectBestModel} method employs a Least-Recently-Used (LRU) style selection. It estimates the required tokens based on character count, applies a safety buffer, and filters available models to ensure they have sufficient context limits. Models are then sorted primarily by their daily \texttt{RequestCount} (to distribute the load evenly) and secondarily by \texttt{ContextLength}.

\textbf{Resilient Rate-Limit Handling:}
When a model responds with \texttt{429 Too Many Requests}, the service does not fail; it parses the specific \texttt{x-ratelimit-reset} header. Because OpenRouter proxies various providers, the service handles multiple time formats:
\begin{itemize}
    \item \textbf{Suffixes:} e.g., \texttt{10s} (seconds), \texttt{5m} (minutes), \texttt{24h} (hours).
    \item \textbf{UNIX Timestamps:} Intelligently distinguishing between seconds and milliseconds based on magnitude.
\end{itemize}
The exhausted model's \texttt{ResetTime} is updated, and the while-loop automatically pivots to the next best model. Furthermore, if a model returns \texttt{402 Payment Required}, indicating a bait-and-switch on "free" pricing, it is permanently excised from the local \texttt{\_models} dictionary.

\subsection{Context-Construction Logic (RAG)}
To prevent hallucinations and guarantee determinism, \texttt{TripPlannerService.cs} implements a sophisticated Retrieval-Augmented Generation (RAG) \cite{rag_lewis} pipeline that tightly bounds the AI's search space.

\textbf{Dual-Layer Caching Architecture:}
Before executing complex geographical queries, the service utilizes a dual-layer cache structure to intercept identical requests:
\begin{enumerate}
    \item \textbf{Level 1 (In-Memory Cache):} Extremely fast lookup via \texttt{ICacheService} (backed by \texttt{IMemoryCache}) using a complex composite key combining coordinates, location name, duration, and user preferences.
    \item \textbf{Level 2 (DB Persistence):} If the ephemeral cache is purged, the \texttt{PrecomputedTripPlan} table in the database acts as a durable fallback. A cache hit here repopulates the L1 cache dynamically.
\end{enumerate}

\textbf{Spatial Filtering \& The 30km Constraint:}
If a cache miss occurs, the service queries the local database to compute point-to-point geographic distances. To ensure the AI doesn't cross city borders (e.g., suggesting an attraction in Alexandria when the user is in Cairo), a strict 30-kilometer maximum radius constraint is enforced. Any location falling outside this geographic boundary is discarded entirely, ensuring the itinerary routing remains physically possible.

\textbf{Strict Local Data Guarantee (Zero On-the-Fly Fetching):}
A major architectural decision was made to entirely eliminate dynamic API fallbacks during generation. We discovered that attempting to fetch missing data from external sources like Google Places on-the-fly introduced unpredictable latency spikes that ruined the user experience. Instead, the system does \textit{not} fallback to fetching Google Places on-the-fly. The RAG pipeline relies completely on the meticulously pre-populated local database to guarantee consistently low latency. By shifting the burden of data aggregation to offline seeding scripts, we achieved a lightning-fast, deterministic generation phase.

\textbf{Categorical Injection and Token Compression:}
To drastically lower the prompt usage and reduce API costs, we implemented a strict token compression strategy. Instead of sending the full JSON payload of all available locations to the AI, we first categorically limit the data pool. The system selects exactly the highest-rated 3 hotels, 20 attractions, 8 restaurants, and 4 cafes. 

Then, rather than sending lengthy descriptions, we compress each selected location into a single dense string format: \texttt{[TYPE] Name | Rating | Categories | Address}. This simple, highly efficient formatting minimizes token consumption while providing the LLM with the exact necessary context to reason about scheduling and proximity, effectively lowering context costs without sacrificing intelligence.

\subsection{System and User Prompt Definitions}
The final prompts merged in May 2026 enforce strict behavioral rules on the LLM.

\textbf{System Prompt (Trip Generation):}
The system prompt instructs the LLM to act as an expert Egyptian travel planner and mandates that the output strictly matches a predefined JSON structure containing daily themes and activities. Furthermore, it enforces rigid rules such as only using the exact place names provided in the context, adhering to realistic time slots, and ensuring a balanced daily mix of hotels, restaurants, attractions, and cafes.

\textbf{User Prompt (Trip Generation):}
The user prompt dynamically injects the user's desired duration, location, budget, interests, and preferences, along with the compressed list of candidate places generated by the RAG pipeline. It explicitly instructs the AI to generate a full itinerary using only the provided entities.

\subsection{The Swapping Mechanism}
Kemora allows users to reject a specific activity and ask the AI for an alternative. This utilizes a separate, highly targeted prompt:

\textbf{Swapping Prompt:}
When a user rejects an activity, the swapping prompt provides the current place and the user's reason for rejection, instructing the AI to return a high-quality alternative from the same area in a strict JSON format.

\subsection{Defensive Parsing and Error Recovery}
Despite rigorous prompt engineering to enforce JSON structure, LLMs occasionally produce malformed outputs---such as truncated JSON arrays, unescaped quotes, or conversational filler (e.g., "Here is your JSON:"). To ensure application stability, the system implements defensive parsing mechanisms using custom deserialization logic that strips markdown code blocks and conversational text before parsing. If deserialization fails, a retry policy is triggered, dynamically reprompting the LLM with the parsing error to self-correct the structure. This error recovery loop dramatically reduces failure rates in production.

\section{Caching and Performance Optimizations}
In late June 2026, profiling under simulated load revealed a critical bottleneck: the backend was attempting to enrich all candidate places sequentially \textit{before} feeding them to the AI. Because Entity Framework Core's \texttt{DbContext} is not thread-safe, this concurrent fetching and saving in a loop led to frequent \texttt{InvalidOperationException} crashes and unacceptable latency. We had to rethink our approach entirely to restore system stability.

\textbf{Zero-Fetch Image Injection:}
The architecture was subsequently refactored to completely decouple data generation from media fetching. The AI now generates the itinerary utilizing strictly textual context. Once the itinerary is parsed, the backend avoids external image API calls entirely. Instead, it systematically scans the generated JSON, matches the AI's chosen locations against our pre-loaded local dataset, and injects the corresponding image URLs directly into the final response object. This transition from sequential bulk fetching to local, in-memory image injection eradicated the API bottleneck, massively improving response time.

\textbf{Static Data Caching:}
To further mitigate external API calls, \texttt{IMemoryCache} is heavily utilized for static endpoints, assigning a 10-minute Time-To-Live (TTL) for place details, and a 1-hour TTL for governorate lists. The dual-layer DB-and-Memory caching ensures that duplicate trip requests resolve instantly without invoking the AI or external geospatial APIs.

\textbf{Quantitative Load Testing Data:}
To empirically validate the impact of our optimizations, we conducted rigorous load testing. Before implementing the dual-layer cache and token compression strategies, average API latency for trip generation exceeded 12 seconds, with significant memory overhead due to extensive concurrent database connections. Post-optimization metrics demonstrated a remarkable improvement: API latency plummeted to under 2.5 seconds for cache misses and less than 100 milliseconds for cache hits. Furthermore, the categorical token compression strategy reduced the average prompt size by 65\%, resulting in a proportional reduction in LLM inference costs while stabilizing memory consumption.

\section{Enterprise Resiliency Patterns}
To ensure robust interactions with external APIs like OpenRouter \cite{openrouter}, the backend integrates the Polly library to implement enterprise-grade resiliency patterns. 

\textbf{Circuit Breaker and Bulkhead Isolation:}
Instead of allowing repeated timeouts to exhaust server resources during an OpenRouter outage or severe rate-limiting event, we implemented a Circuit Breaker policy using Polly. If a consecutive number of requests fail, the circuit opens, instantly failing subsequent requests to prevent cascading failures and allow the external service time to recover. Concurrently, a Bulkhead Isolation policy limits the number of concurrent connections to the AI service. This guarantees that a spike in trip generation requests will not saturate the thread pool, ensuring that other application components---such as authentication and user profile management---remain highly available even when external AI providers degrade.

\section{Conflict Resolution Strategy for Offline-First Sync}
To support users traveling in areas with poor cellular connectivity, the mobile client employs an offline-first architecture. Users can interact with and modify their saved trip plans without an active internet connection, with changes cached locally.

\textbf{Synchronization and Conflict Policies:}
When cellular connectivity is regained, a background synchronization process reconciles local modifications with the remote database. To handle concurrent edits from multiple devices, we established a strict conflict resolution strategy. The system primarily relies on a Timestamp-based policy, utilizing an \texttt{UpdatedAt} timestamp to determine the most recent mutation. In complex merging scenarios where precise timestamps cannot reliably resolve the conflict, the system falls back to a Server-Wins policy, prioritizing the authoritative state on the backend to maintain overall data integrity. This robust synchronization ensures a seamless user experience despite fluctuating network conditions.

\section{DevOps and Production Monitoring}
To maintain high availability and streamline deployment, the Kemora backend was containerized using Docker, allowing for consistent environments across development and production. For comprehensive observability, we integrated Serilog and Application Insights. This monitoring setup captures structured logs, traces incoming HTTP requests, and visualizes AI generation times. By establishing targeted dashboards, the engineering team can proactively identify rate-limit events, track model response degradation, and monitor overall system health in real-time, ensuring the optimized setup remains robust under production load.
